% ============================================================================
\clearpage
\section{Transporte, dinámica y extracción de señal}\label{sec:transporte}
% ============================================================================

\ficha{Distribuciones $\Q_t$ a plazo constante con bandas y soporte identificado; información base de
precio, volatilidad y mercado.}
{Representar por cuantiles, medir distancias y cambios firmados, reducir dimensión, filtrar causalmente
y calcular la innovación.}
{Series de distancias, factores filtrados con incertidumbre, variables de cola e innovaciones.}
{Todo ajuste usa solo el pasado; reconstrucciones ordenadas; medidas recortadas declaradas como tales.}

\subsection{Representación por cuantiles relativos al forward}\label{sec:cuantiles}

Cada distribución se representa por su función cuantil de rendimientos logarítmicos relativos al
forward, a plazo constante $\tau$:
\begin{equation}
x_t(u,\tau)=F^{-1}_{\Q_t}(u),\qquad \text{con } x=\ln\!\bigl(S_{t+\tau}/F_{t,\tau}\bigr),\quad u\in(0,1).
\label{eq:cuantil}
\end{equation}
Medir respecto del forward retira el desplazamiento puramente mecánico del nivel de precios: si el
subyacente sube 1\% y nada más cambia, $x_t(u)$ no se mueve. La vista en precios absolutos se conserva
por separado para interpretar la posición. El plazo constante se obtiene interpolando $\theta_T$ entre
vencimientos listados; como $\theta$ interpolado linealmente sigue siendo no decreciente, se preserva la
ausencia de arbitraje de calendario. Por construcción $\E^{\Q}\bigl[e^{x}\bigr]=1$.

\subsection{Distancias entre distribuciones y cambios firmados}\label{sec:distancias}

En una dimensión, el transporte óptimo entre dos distribuciones es el \emph{reordenamiento monótono}: cada
cuantil se desplaza a su homólogo, $T=F_t^{-1}\circ F_{t-1}$ \citep{peyre2019}. La distancia de
Wasserstein de orden 2 se calcula entonces directamente sobre las funciones cuantil:
\begin{equation}
W_2^2(\Q_t,\Q_{t-1})=\int_0^1\bigl[x_t(u)-x_{t-1}(u)\bigr]^2\,du,\qquad
W_1(\Q_t,\Q_{t-1})=\int_0^1\bigl|x_t(u)-x_{t-1}(u)\bigr|\,du .
\label{eq:w2}
\end{equation}
$W_1$ es menos sensible que $W_2$ a desplazamientos extremos y equivale al área entre las dos funciones
de distribución. Para familias de posición y escala, $W_2^2=(\Delta\mu)^2+(\Delta\sigma)^2$, lo que ayuda
a leer su magnitud. Si solo existe cobertura fiable del 5\% al 95\%, se reporta la versión recortada
$W_{2,[a,b]}^2=\int_a^b[\Delta x(u)]^2du$ con su rango, nunca como distancia de toda la distribución.

\grafica[0.6]{s3_densidades_choque}{Distribución implícita a 30 días antes y después de un choque común en
el mundo sintético: más dispersión y una cola izquierda más larga.}{fig:densidades_choque}

\grafica[0.6]{s3_mapa_transporte}{Mapa de transporte óptimo entre las dos sesiones de la
\figref{fig:densidades_choque}: cada cuantil anterior (eje horizontal) se envía al mismo cuantil posterior
(eje vertical). La pendiente mayor que uno indica más dispersión; la diagonal sería ``sin cambio''.}{fig:mapa_transporte}

\grafica[0.6]{s3_funciones_cuantil}{Funciones cuantil antes y después del choque. $W_2$ es la distancia
$L^2$ entre ellas; la leyenda compara la versión completa, la recortada a $[0.05,0.95]$ y $W_1$.}{fig:funciones_cuantil}

La magnitud sola no dice qué cambió. Por eso se registran también los \emph{cambios firmados}
$\Delta x_t(u)$ en percentiles clave, que distinguen deformaciones con distancias parecidas
(figuras~\ref{fig:cambios_firmados} y~\ref{fig:distancias_arquetipos}).

\grafica[0.6]{s3_cambios_firmados}{Perfil firmado de tres deformaciones arquetípicas de una misma
distribución: más dispersión, cola izquierda más pesada y cambio solo en las alas. Un $\Delta x<0$ con
$u<1/2$ significa que la cola izquierda se alarga.}{fig:cambios_firmados}

\grafica[0.6]{s3_distancias_arquetipos}{Distancias de los tres arquetipos. $W_2$ pesa más el cambio en las
alas que $W_1$; la versión recortada casi no lo ve. Ninguna cifra aislada identifica el tipo de cambio.}{fig:distancias_arquetipos}

\grafica[0.64]{s3_mapa_cuantiles}{Mapa temporal de cuantiles: desviación de cada $x_t(u)$ respecto de su
media de entrenamiento, por sesión. El choque común de la sesión 170 alarga ambas colas; las columnas en
blanco son snapshots faltantes.}{fig:mapa_cuantiles}

\grafica[0.64]{s3_distancias_diarias}{Distancias entre sesiones consecutivas. Los puntos rojos son sesiones
con spreads amplios, donde parte de la distancia es ruido de ajuste y no cambio de distribución: una
razón para filtrar antes de interpretar.}{fig:distancias_diarias}

\subsection{Cadencia, velocidad y aceleración}\label{sec:cadencia}

Las velocidades se definen sobre una cadencia fija y con diferencias hacia atrás,
$v_t(u)=x_t(u)-x_{t-1}(u)$, en canales separados: \emph{cierre--apertura} y \emph{apertura--apertura}.
Fines de semana, feriados y ventanas intradía son intervalos de tipo distinto y se etiquetan como tales;
no se mezclan como si fueran idénticos ni se escalan por tiempo calendario sin un modelo. La aceleración
$a_t(u)=v_t(u)-v_{t-1}(u)$ queda como candidata secundaria por una razón aritmética: si cada $x_t$ tiene
ruido de medición independiente de varianza $\sigma^2$, la velocidad hereda $2\sigma^2$ y la aceleración
$6\sigma^2$ (figuras~\ref{fig:senal_ruido} y~\ref{fig:aceleracion}).

\grafica[0.6]{s3_senal_ruido}{Cociente entre la varianza de la señal y la del ruido de medición para el
nivel, la velocidad y la aceleración de un factor sintético. Diferenciar multiplica la desviación estándar
del ruido por $\sqrt2$ y $\sqrt6$ mientras la señal se vuelve más pequeña.}{fig:senal_ruido}

\grafica[0.6]{s3_aceleracion}{Aceleración latente frente al ruido de medición que contiene la aceleración
observada: fuera del choque son de magnitud comparable.}{fig:aceleracion}

\subsection{Reducción de dimensión: componentes funcionales y variables de cola}\label{sec:fpca}

Para resumir la forma se aplica un análisis de componentes principales funcional a los cuantiles
centrados, con pesos de cuadratura en la malla de $u$ y ajustado \emph{únicamente con datos de
entrenamiento}:
\begin{equation}
x_t(u)\approx\bar x(u)+\sum_{j=1}^{J}s_{tj}\,\phi_j(u),\qquad
s_{tj}=\int_0^1\bigl[x_t(u)-\bar x(u)\bigr]\phi_j(u)\,du .
\label{eq:fpca}
\end{equation}
Se conservan factores suficientes para explicar dispersión y asimetría. Pero una PCA global puede ignorar
las colas precisamente porque explican poca varianza total (figuras~\ref{fig:fpca_autofunciones}
y~\ref{fig:fpca_varianza}); por eso se añaden variables explícitas de cola, por ejemplo el ancho de la cola
izquierda $x_t(0.10)-x_t(0.01)$ dentro del rango identificado, o la asimetría por cuantiles
$[x_t(0.9)+x_t(0.1)-2x_t(0.5)]/[x_t(0.9)-x_t(0.1)]$.

\grafica[0.6]{s3_fpca_autofunciones}{Primeras autofunciones de la PCA funcional, estimadas solo con las
sesiones de entrenamiento del mundo sintético. La primera captura dispersión; la segunda, asimetría; la
tercera, las alas.}{fig:fpca_autofunciones}

\grafica[0.6]{s3_fpca_varianza}{Varianza explicada por componente, en escala logarítmica. La componente
de colas explica una fracción ínfima de la varianza total aunque pueda ser la más relevante para el
riesgo.}{fig:fpca_varianza}

Toda reconstrucción debe seguir siendo una función cuantil válida, es decir, creciente en $u$. Con pocos
factores eso no está garantizado; se corrige con el reordenamiento creciente \citep{chernozhukov2010}, que
no aumenta la distancia $L^p$ a la función verdadera, o se trabaja en una transformación que garantice
monotonía, como el logaritmo de la densidad cuantílica \citep{petersen2016}.

\subsection{Filtrado causal con incertidumbre de observación}\label{sec:kalman}

Los factores observados cada sesión mezclan cambio real y error de ajuste. Se filtran con un modelo de
estado lineal y el filtro de Kalman \citep{kalman1960}:
\begin{equation}
f_t=A f_{t-1}+\eta_t,\quad \eta_t\sim\mathcal N(0,Q);\qquad
y_t=f_t+\varepsilon_t,\quad \varepsilon_t\sim\mathcal N(0,R_t),
\label{eq:kalman}
\end{equation}
donde la varianza de observación $R_t$ \emph{depende de la calidad} de la sesión: ancho de spreads, número
de strikes válidos y error de ajuste. Con calidad baja, el filtro confía menos en la observación y más en
la dinámica. Si falta el snapshot, solo se ejecuta la predicción. Los parámetros $(A,Q,R_0)$ se estiman
por máxima verosimilitud con datos de entrenamiento.

Solo se usa el \emph{filtro}, que emplea información hasta $t$. Un suavizador que incorpore observaciones
futuras produce series más limpias y backtests inválidos. La comparación obligada es con una media
exponencial (EWMA) y con la persistencia: un filtro más complejo puede borrar señales rápidas
(figuras~\ref{fig:kalman} y~\ref{fig:kalman_error}).

\grafica[0.64]{s3_kalman}{Filtrado del primer factor en el mundo sintético. El filtro de Kalman, con $R_t$
mayor en las sesiones de spreads amplios, sigue al factor latente; la EWMA suaviza el ruido pero llega
tarde al choque de la sesión 170.}{fig:kalman}

\grafica[0.6]{s3_kalman_error}{Error frente al factor latente por ventana de evaluación. Ningún método
domina: la persistencia es casi tan buena con calidad normal, el filtro ayuda cuando los spreads son
amplios y la EWMA falla tras el choque.}{fig:kalman_error}

\subsection{La propuesta diferenciadora: la innovación}\label{sec:innovacion}

Buena parte del cambio diario de la distribución implícita es predecible con información que ya está en
el precio: si el mercado cae, la cola izquierda suele alargarse. Lo potencialmente informativo es lo que
\emph{no} se explica así. Se define la innovación
\begin{equation}
\epsilon_t(u)=\Delta x_t(u)-\widehat{\E}\bigl[\Delta x_t(u)\mid\mathcal{I}_t^{base}\bigr],
\label{eq:innovacion}
\end{equation}
donde $\mathcal I_t^{base}$ incluye el retorno ya observado del subyacente, su volatilidad, el movimiento
de mercado y sector, el gap de apertura y el calendario conocido de resultados y dividendos
(\figref{fig:innovacion_diagrama}). El modelo de expectativa se entrena \emph{solo con fechas anteriores};
para una acción, además, se calcula la innovación relativa al mercado,
$\epsilon^A_t-\beta\,\epsilon^M_t$, para separar señal idiosincrática de choque común.

\begin{figure}[H]
\centering
\resizebox{\linewidth}{!}{% Construcción de la innovación de la distribución implícita.
\begin{tikzpicture}[
  n/.style={caja=#1,text width=3.1cm,minimum height=13mm},
  node distance=6mm and 7mm]
\node[n=qazul] (obs) at (0,0) {\textbf{Cambio observado}\\ $\Delta x_t(u)=x_t(u)-x_{t-1}(u)$};
\node[n=qnaranja] (base) at (0,-2.3) {\textbf{Información base} $\mathcal{I}^{base}_t$\\ retorno ya observado,
  volatilidad, mercado y sector, gap, calendario};
\node[n=tenue] (mod) at (4.1,-2.3) {\textbf{Modelo de expectativa}\\ ajustado solo con fechas\\ anteriores a $t$};
\node[n=tenue] (esp) at (8.2,-2.3) {\textbf{Parte esperada}\\ $\widehat{\E}\bigl[\Delta x_t(u)\mid\mathcal{I}^{base}_t\bigr]$};
\node[circle,draw=tinta2,line width=0.7pt,minimum size=6mm,font=\sffamily\bfseries] (resta) at (8.2,0) {$-$};
\node[n=qazul] (inn) at (12.2,0) {\textbf{Innovación}\\ $\epsilon_t(u)$};
\node[n=qnaranja] (uso) at (12.2,-2.3) {\textbf{Variable candidata}\\ para pronosticar $\P$\\ después de $t$};
\draw[flecha] (obs) -- (resta);
\draw[flecha] (base) -- (mod);
\draw[flecha] (mod) -- (esp);
\draw[flecha] (esp) -- (resta);
\draw[flecha] (resta) -- (inn);
\draw[flecha] (inn) -- (uso);
\node[nota,anchor=north] at (6.1,-3.25) {La innovación no prueba información privada ni causalidad: es un candidato cuyo valor
  predictivo se mide fuera de muestra.};
\end{tikzpicture}
}
\caption{Construcción de la innovación. La resta se hace con una expectativa estimada exclusivamente con
el pasado.}
\label{fig:innovacion_diagrama}
\end{figure}

\grafica[0.6]{s3_innovacion_dispersion}{Cambio observado del cuantil 5\% de la acción A frente al esperado
dado la información base. La distancia vertical a la diagonal es la innovación. El evento idiosincrático
(círculo) alarga la cola sin movimiento de precio; el choque común (cuadrado) se explica solo en
parte.}{fig:innovacion_dispersion}

\grafica[0.64]{s3_innovacion_serie}{Innovación del mercado e innovación de la acción A relativa al mercado.
En la versión relativa el choque común casi se cancela y el evento propio de A sobresale.}{fig:innovacion_serie}

\begin{noconcluir}
La innovación no prueba información privada ni causalidad. Es una variable candidata para pronosticar
rendimientos \emph{posteriores} al instante de decisión, y su utilidad se mide fuera de muestra con el
protocolo de la \secref{sec:evidencia}. En el mundo sintético de estas figuras, por diseño, la innovación
no predice nada: solo se muestra cómo se calcula.
\end{noconcluir}

\subsection{Experimentos concretos}\label{sec:experimentos}

La \tabref{tab:experimentos} lista los primeros experimentos. En ninguno se asume una dirección rentable
antes de medirla.

\begin{table}[H]
\centering\small
\caption{Primeros experimentos sobre la señal.}
\label{tab:experimentos}
\begin{tabularx}{\linewidth}{@{}>{\RaggedRight}p{3.4cm} Y Y@{}}
\toprule
\encab{Experimento} & \encab{Variable} & \encab{Pregunta}\\
\midrule
Cola izquierda frente al mercado & $\epsilon^A_t(u)-\beta\,\epsilon^M_t(u)$ para $u\le 0.10$ & ¿Aporta la
deformación idiosincrática de la cola a los pronósticos de riesgo o rendimiento?\\
Deformación entre plazos & Diferencia de $\Delta x_t(u)/\sqrt{\tau}$ entre 7, 30 y 60 días & ¿Distingue
eventos de corto plazo de choques de nivel? Requiere cotizaciones suficientes en cada plazo.\\
Persistencia & Autocorrelación de $\epsilon_t$ y su vida media & ¿La innovación se revierte, persiste o se
incorpora al precio?\\
Idiosincrático frente a común & Descomposición de $\epsilon^A_t$ en componente de mercado y residual & ¿Qué
parte, si alguna, contiene información que no está en el mercado?\\
\bottomrule
\end{tabularx}
\end{table}

\grafica[0.6]{s3_plazos_comun}{Cambio de cuantiles por plazo, normalizado por $\sqrt{\tau}$, ante un
choque común de volatilidad: el perfil es parecido en los tres plazos, algo mayor en el corto.}{fig:plazos_comun}

\grafica[0.6]{s3_plazos_evento}{Ante un evento dentro de los próximos 7 días, la distribución a 7 días se
deforma mucho más que las de 30 y 60: la diferencia entre plazos es la huella del
evento.}{fig:plazos_evento}

\subsection{Dinámica de distribuciones como objetos}\label{sec:war}

Como modelo avanzado se compara la autorregresión de factores con una autorregresión en espacio de
Wasserstein \citep{zhang2022}. La idea: el espacio de distribuciones no es lineal, pero alrededor del
baricentro de Wasserstein $\bar Q$ tiene un espacio tangente que sí lo es. En una dimensión el baricentro
es el promedio de funciones cuantil, $\bar x(u)=\E[x_t(u)]$, y el ``logaritmo'' de $Q_t$ es simplemente
$v_t(u)=x_t(u)-\bar x(u)$. Se ajusta un autorregresivo sobre $v_t$ y se regresa a funciones cuantil,
proyectando cuando el pronóstico deja de ser creciente (\figref{fig:tangente}).

\begin{figure}[H]
\centering
\resizebox{0.86\linewidth}{!}{% Autorregresión en el espacio tangente del baricentro de Wasserstein.
\begin{tikzpicture}[x=1cm,y=1cm]
% Superficie curva: espacio de distribuciones con la métrica W2
\fill[qazul!6] (0,0) .. controls (2.5,1.6) and (6.5,1.6) .. (9,0.2)
  .. controls (9.6,-1.2) and (9.2,-2.6) .. (8.4,-3.4)
  .. controls (6,-2.1) and (2.8,-2.2) .. (0.4,-3.2)
  .. controls (-0.3,-2.2) and (-0.4,-1) .. (0,0);
\draw[qazul!60,line width=0.6pt] (0,0) .. controls (2.5,1.6) and (6.5,1.6) .. (9,0.2)
  .. controls (9.6,-1.2) and (9.2,-2.6) .. (8.4,-3.4)
  .. controls (6,-2.1) and (2.8,-2.2) .. (0.4,-3.2)
  .. controls (-0.3,-2.2) and (-0.4,-1) .. (0,0);
\node[nota,anchor=west] at (0.3,-2.75) {distribuciones con la métrica $W_2$ (no lineal)};
% Plano tangente en el baricentro
\coordinate (B) at (4.6,-0.9);
\fill[qnaranja!10,draw=qnaranja!70,line width=0.6pt] (1.9,0.35) -- (7.9,0.35) -- (7.2,-2.05) -- (1.2,-2.05) -- cycle;
\node[nota,anchor=east,text=tinta] at (7.05,-1.8) {espacio tangente $T_{\bar{Q}}$ (lineal)};
\fill[tinta] (B) circle (2.2pt);
\node[font=\sffamily\scriptsize,anchor=north] at (B) {$\bar{Q}$};
% Log de dos sesiones y predicción
\draw[flecha,draw=qazuloscuro] (B) -- (2.6,-0.3) node[anchor=east,font=\sffamily\scriptsize]{$v_{t-1}$};
\draw[flecha,draw=qazuloscuro] (B) -- (3.4,0.05) node[anchor=south,font=\sffamily\scriptsize]{$v_t$};
\draw[flecha,draw=qnaranja,line width=0.9pt] (B) -- (5.4,0.1) node[anchor=south west,font=\sffamily\scriptsize]{$\hat v_{t+1}=A\,v_t$};
% Distribuciones en la superficie
\fill[qazul] (1.7,1.05) circle (2pt) node[anchor=south,font=\sffamily\scriptsize]{$Q_{t-1}$};
\fill[qazul] (3.2,1.25) circle (2pt) node[anchor=south,font=\sffamily\scriptsize]{$Q_t$};
\fill[qnaranja] (6.3,1.2) circle (2pt) node[anchor=south,font=\sffamily\scriptsize]{$\widehat{Q}_{t+1}$};
\draw[flecha,draw=tenue,densely dotted] (1.7,1.0) -- (2.55,-0.2);
\draw[flecha,draw=tenue,densely dotted] (3.2,1.2) -- (3.4,0.15);
\draw[flecha,draw=qnaranja,densely dotted] (5.45,0.2) -- (6.25,1.1);
\node[nota,anchor=west,align=left] at (9.4,1.0) {$\exp$ y proyección a\\ funciones cuantil crecientes};
\node[nota,anchor=east,align=right] at (-0.45,0.6) {$\log$ en una dimensión:\\ $v_t(u)=x_t(u)-\bar{x}(u)$};
\end{tikzpicture}
}
\caption{Autorregresión en el espacio tangente del baricentro. En una dimensión el espacio tangente
coincide con el de cuantiles centrados, así que la PCA funcional de la \secref{sec:fpca} y este modelo son
parientes cercanos: difieren en cómo parametrizan la dinámica.}
\label{fig:tangente}
\end{figure}

Esta familia permite pronosticar distribuciones como objetos. Su uso aquí es una adaptación de
investigación: no demuestra capacidad para anticipar el precio de las acciones.

\subsection{Límites del transporte}\label{sec:limites}

El transporte óptimo entre marginales \emph{no} identifica un proceso temporal del precio ni probabilidades
de tocar un stop. Las figuras~\ref{fig:trayectorias_difusion} y~\ref{fig:trayectorias_transporte} lo
muestran con 20\,000 trayectorias: ambas tienen exactamente las mismas distribuciones a 30 y 60 días, pero
en una se reordenan como si cada cuantil siguiera su propio camino.

\grafica[0.6]{s3_trayectorias_difusion}{Trayectorias de una difusión (muestra de 70). La probabilidad de
tocar un stop en 90 antes de 60 días es la que indica el título.}{fig:trayectorias_difusion}

\grafica[0.6]{s3_trayectorias_transporte}{Las mismas marginales leídas como trayectorias del mapa monótono:
la probabilidad de tocar el stop cae casi a la mitad. Las marginales no deciden cuál de las dos es la
correcta.}{fig:trayectorias_transporte}

\begin{ideaclave}
El campo de transporte es una descripción geométrica de cómo cambió la distribución, no una ley física
que deba imponerse al mercado. Las probabilidades de trayectoria (stops, barreras, máximos) requieren un
modelo dinámico explícito y su propia validación.
\end{ideaclave}

\begin{criterio}
Separar ruido de cambios persistentes con filtros causales; todas las series de esta etapa calculables sin
información futura y reproducibles con el snapshot de cada sesión.
\end{criterio}
